\clearpage

\section{基于RT-DETRv2的RHEED条纹与斑点检测算法}

本章节围绕RHEED图像中的条纹与斑点目标检测任务展开研究。考虑到论文的工程化落地与可复现性要求，本文采用\textbf{RT-DETRv2}作为主方法完成检测实验。RT-DETRv2是一类端到端的Transformer检测器，具有推理速度快、收敛稳定、易于迁移的特点，适合在小样本工业数据上进行快速微调。本文首先介绍RT-DETRv2的检测流程与核心原理；随后给出数据标注与训练流程；最后在实验结果中对检测性能进行分析，并将DiffusionDet作为未来拓展工作进行讨论。

\subsection{流程介绍}

RT-DETRv2沿用DETR的集合预测框架，使用卷积主干网络提取特征并输入Transformer编码器与解码器，通过一组固定数量的查询向量（queries）完成目标定位与类别预测。其核心思想是将检测任务转化为集合匹配问题，通过匈牙利算法完成预测框与真实框的一一匹配，从而避免传统检测器中的复杂候选框筛选过程。模型输出的检测集合记为：
\begin{equation}
    \hat{y} = \left\{ (\hat{p}_i, \hat{b}_i) \right\}_{i=1}^{N}
\end{equation}
其中$\hat{p}_i$表示第$i$个预测框的类别概率，$\hat{b}_i=(c_x, c_y, w, h)$表示对应的边界框参数。模型训练时采用集合匹配损失：
\begin{equation}
    \mathcal{L} = \lambda_{\text{cls}} \mathcal{L}_{\text{cls}} + \lambda_{\text{l1}} \mathcal{L}_{\text{l1}} + \lambda_{\text{giou}} \mathcal{L}_{\text{giou}}
\end{equation}
其中$\mathcal{L}_{\text{cls}}$为分类损失，$\mathcal{L}_{\text{l1}}$和$\mathcal{L}_{\text{giou}}$用于优化框回归与几何一致性。相较于传统两阶段检测器，RT-DETRv2通过端到端集合优化避免了NMS等后处理步骤，更适合部署到实际外延过程的自动监控系统中。

针对RHEED图像的特点（背景相对单一、目标结构稀疏、尺度变化明显），RT-DETRv2具备以下匹配优势：
\begin{itemize}
    \item \textbf{全局建模能力强}：Transformer编码器能够捕捉条纹与斑点之间的全局空间关系，有利于抑制噪声背景。
    \item \textbf{小样本微调稳定}：可直接加载COCO预训练权重，在少量标注数据上快速收敛。
    \item \textbf{端到端结构简洁}：训练与推理流程统一，易于工程复现与部署。
\end{itemize}

\subsection{易复现性与工程实现}

本文检测实验采用公开实现的RT-DETRv2版本，加载COCO预训练权重，在完成RHEED条纹/斑点标注后进行微调。整体复现实验仅需完成\textbf{数据标注、数据拆分、训练与结果可视化}四个步骤，不依赖复杂的多阶段预训练流程。在单卡A100环境下，RT-DETRv2能够在数小时内完成收敛（具体时间与数据规模相关），适合在10天内完成实验闭环并形成论文结果。

\subsection{DiffusionDet作为未来拓展工作}

DiffusionDet通过扩散去噪过程逐步生成目标检测框，天然适用于高噪声、小目标、密集目标等复杂场景。在RHEED图像中，条纹与斑点存在尺度变化与局部遮挡，扩散模型的逐步去噪机制有望提升鲁棒性与定位精度，尤其适用于高密度量子点生长过程的结构检测。然而该方法依赖detectron2体系，训练流程复杂、复现成本高，因此本文将其作为未来拓展方向，待标注规模扩大后进一步验证其潜力。

\subsection{实验设计}

\textbf{数据集构建}：在\texttt{rheed\_images\_png}中抽取样本进行人工标注，类别设置为\texttt{stripe}与\texttt{spot}，采用统一的边界框标注规范，并划分训练集/验证集/测试集。

\textbf{数据预处理}：对图像进行统一尺寸缩放，保证条纹与斑点在不同批次图像中尺度一致；训练阶段使用随机翻转与轻度亮度扰动，以增强模型对噪声的鲁棒性。

\textbf{训练策略}：以COCO预训练权重为初始化，采用端到端微调策略，训练超参数保持简洁可控，避免过度调参。

\textbf{评价指标}：使用mAP$_{50}$、mAP$_{50:95}$、召回率（Recall）与检测可视化结果作为评价标准，并统计典型失败案例进行误检分析。

\subsection{实验结果分析}

为保证实验可复现性与论文结果可追溯性，本节给出三种从头预训练检测模型（RT-DETR、RT-DETRv2、D-FINE）的关键观测值、验证集结果与测试集结果。对应数据来源如下：
\begin{itemize}
    \item 训练/验证日志：\texttt{outputs/rtdetr\_r18vd\_50e\_rheed/log.txt}、\texttt{outputs/rtdetrv2\_r18vd\_50e\_rheed/log.txt}、\texttt{outputs/dfine\_rheed/log.txt}
    \item 测试评估产物：\texttt{outputs/rtdetr\_r18vd\_50e\_rheed\_test/eval.pth}、\texttt{outputs/rtdetrv2\_r18vd\_50e\_rheed\_test/eval.pth}、\texttt{outputs/dfine\_rheed\_test/eval.pth}
\end{itemize}

\subsubsection{从头预训练关键观测值（用于可视化）}

表\ref{tab:cp3_obs_metrics_by_model}给出了三种模型在训练早期（epoch 0）、中期（epoch 25）与末期（epoch 49）的关键观测值。可据此绘制\texttt{train\_loss-epoch}与\texttt{AP-epoch}曲线，直观比较收敛速度与性能上限。

\begin{table}[h]
    \centering
    \caption{从头预训练关键观测值对比（指标为行，模型为列；除train\_loss外均为\%）}
    \label{tab:cp3_obs_metrics_by_model}
    \begin{tabular}{lccc}
        \hline
        \textbf{指标}            & \textbf{RT-DETR} & \textbf{RT-DETRv2} & \textbf{D-FINE} \\ \hline
        train\_loss@epoch0     & 14.18            & 20.05              & 28.46           \\
        train\_loss@epoch25    & 5.17             & 7.73               & 15.68           \\
        train\_loss@epoch49    & 4.39             & 5.79               & 13.23           \\
        AP@[0.50:0.95]@epoch0  & 28.32\%          & 0.00\%             & 0.00\%          \\
        AP@[0.50:0.95]@epoch25 & 91.09\%          & 82.40\%            & 93.71\%         \\
        AP@[0.50:0.95]@epoch49 & 91.56\%          & 90.56\%            & 98.67\%         \\ \hline
    \end{tabular}
\end{table}

\subsubsection{验证集检测结果}

验证集结果由各模型训练日志中\texttt{epoch=49}对应的\texttt{test\_coco\_eval\_bbox}提取，指标采用COCO标准定义。结果见表\ref{tab:cp3_val_metrics_by_model}。

\begin{table}[h]
    \centering
    \caption{验证集检测结果对比（指标为行，模型为列）}
    \label{tab:cp3_val_metrics_by_model}
    \begin{tabular}{lccc}
        \hline
        \textbf{指标}    & \textbf{RT-DETR} & \textbf{RT-DETRv2} & \textbf{D-FINE} \\ \hline
        AP@[0.50:0.95] & 0.915630         & 0.905629           & 0.986695        \\
        AP50           & 0.994758         & 0.992264           & 0.994922        \\
        AP75           & 0.984655         & 0.987434           & 0.989939        \\
        AP\_small      & 0.883276         & 0.888985           & 0.987743        \\
        AP\_medium     & 0.948091         & 0.922459           & 0.985737        \\
        AP\_large      & -1.000000        & -1.000000          & -1.000000       \\
        AR@100         & 0.931593         & 0.939622           & 0.992971        \\ \hline
    \end{tabular}
\end{table}

\subsubsection{测试集检测结果（rough统计）}

当前测试集结果来自\texttt{eval.pth}的粗粒度统计（\texttt{precision\_valid\_mean}、\texttt{recall\_valid\_mean}、\texttt{ap50\_mean\_rough}），可用于横向比较，但不等价于\texttt{COCOeval.summarize()}输出的标准AP统计。结果见表\ref{tab:cp3_test_rough_metrics_by_model}。

\begin{table}[h]
    \centering
    \caption{测试集粗统计结果对比（指标为行，模型为列；单位：\%）}
    \label{tab:cp3_test_rough_metrics_by_model}
    \begin{tabular}{lccc}
        \hline
        \textbf{指标}            & \textbf{RT-DETR} & \textbf{RT-DETRv2} & \textbf{D-FINE} \\ \hline
        precision\_valid\_mean & 65.44\%          & 69.14\%            & 70.02\%         \\
        recall\_valid\_mean    & 66.59\%          & 69.49\%            & 70.16\%         \\
        ap50\_mean\_rough      & 99.94\%          & 99.98\%            & 99.99\%         \\ \hline
    \end{tabular}
\end{table}

由上述结果可见，D-FINE在当前三模型对比中取得了最高的验证集指标与测试粗统计指标。后续在第二组实验中，将进一步在RT-DETRv2框架下引入多任务特征共享机制，验证共享分类骨干对检测性能提升的有效性。

\subsection{本章小结}

本章采用RT-DETRv2完成RHEED条纹与斑点检测任务的工程化验证，形成了完整的标注—训练—评估流程，保证了实验的可复现性。同时讨论了DiffusionDet在复杂噪声场景中的潜在优势，并将其作为后续扩展方向。整体而言，本章方法满足小样本工业数据场景下的可部署性与可控复现要求，为外延过程自动监控提供了关键支撑。
